\chapter{Problem Statement, Objectives and Scope}
\label{ch:problem}

\section{Context and Background}
Despite major progress in generative AI and large language models
over the past five years, current story-generation systems remain
fundamentally limited in their ability to create truly personalised
and engaging narrative experiences. The dominant paradigm of
prompt-based generation, while producing impressive individual
outputs, fails to capture the iterative, collaborative nature of
effective storytelling.

Existing models are unable to convert step-by-step user responses
into coherent, personalised and context-rich stories. The single-shot
generation approach requires users to articulate all preferences
upfront, a cognitively demanding task that many users find unnatural
and limiting. This approach also prevents the discovery of preferences
that emerge through interaction --- users often do not know what they
want until they see initial results and can react to them.

\section{Technical Challenges}
Current storytelling systems face several interconnected technical
challenges:

\begin{itemize}
    \item \textbf{Context-window limitations.} LLMs have finite
    context windows that constrain their ability to maintain
    awareness of all relevant story elements across extended
    narratives. As stories grow longer, early elements may be
    forgotten or contradicted.
    \item \textbf{Coherence degradation.} Systems often lose
    narrative consistency in plot, characters and emotional flow when
    multiple user inputs are involved. Each new input can introduce
    inconsistencies that compound over time.
    \item \textbf{Intent interpretation.} Accurately interpreting
    user intent from natural-language responses requires sophisticated
    understanding of context, implication and preference inference
    that current systems handle imperfectly.
    \item \textbf{Balancing agency and quality.} Giving users
    complete control over narrative direction can lead to incoherent
    or poorly structured stories, while limiting control reduces
    personalisation. Finding the right balance is challenging.
    \item \textbf{Multimodal latency and rate limiting.} Augmenting
    text with images, audio and ambient sound introduces additional
    latency and rate-limit failure modes that must be hidden behind a
    smooth UI.
\end{itemize}

\section{User Experience Limitations}
From a user experience perspective, current systems present several
significant limitations:

\begin{itemize}
    \item No framework exists that allows users to control the
    story's theme, tone, genre or direction through an interactive
    Q\&A process.
    \item Generic story outputs fail to reflect individual user
    personalities, interests and emotional states.
    \item The lack of iterative refinement means users cannot guide
    stories toward desired outcomes based on how narratives develop.
    \item Absence of explanation or transparency in generation
    decisions prevents users from understanding why stories develop
    in particular directions.
    \item Existing tools rarely surface reader analytics (such as
    chapter trees or character maps) that would help readers
    navigate or revisit complex narratives.
\end{itemize}

\section{Problem Statement}
Existing AI-based story-generation systems cannot capture user
responses through structured, real-time conversational interaction
and convert them into coherent, multimodal, personalised stories.
There is a clear need for an intelligent system that elicits user
preferences through a guided Q\&A protocol, maintains narrative
coherence across extended interactions, integrates illustration,
narration and ambient music as first-class modalities, and admits
end-user extension of the elicitation protocol --- all while
remaining provider-agnostic, deployable on commodity hardware and
freely accessible.

\section{Research Objectives}

\subsection{Primary Objectives}
\textbf{O1. Structured Input Design.} Design and implement a
comprehensive Q\&A framework that systematically captures user
preferences through natural conversational interaction. The framework
shall support 4, 8, 12 and 17-question budgets, phase-ordered across
setting, character, theme and plot dimensions, and shall offer a
per-question custom-answer override as well as user-added question
authoring.

\textbf{O2. Personalised Story Generation.} Create LLM-based
generation capabilities that produce deeply personalised narratives
reflecting individual user characteristics. The generator shall be
provider-agnostic across Ollama, Groq, Hugging Face and OpenAI and
shall support runtime switching without restart.

\textbf{O3. Narrative Coherence Control.} Ensure generated narratives
maintain consistency and quality across extended multi-turn
interactions. A five-phase story-flow model with word-count thresholds
shall drive Freytag-arc progression; an active consistency checker
shall flag character contradictions, plot-thread abandonment, pacing
issues, setting inconsistencies and tone shifts.

\textbf{O4. Multimodal Augmentation.} Add per-segment scene
illustration, voice narration and genre-conditioned ambient music as
orchestrated companions of the text, dynamically conditioned on the
current segment content and detected genre.

\textbf{O5. Reader Analytics.} Surface a Story Map (branch tree) and a
Character Relationship Graph for readers to navigate the generated
content, with an explicit co-appearance disclaimer in the UI.

\textbf{O6. Story Quality Evaluation.} Establish robust metrics and
evaluation methodologies that combine LLM-based discourse scoring with
structured human Likert ratings across coherence, creativity,
engagement, consistency and personalisation dimensions.

\subsection{Secondary Objectives}
\begin{itemize}
    \item Create reusable, modular components that can be adapted
    for different storytelling domains (education, entertainment,
    therapy, marketing).
    \item Develop well-documented REST API interfaces that enable
    integration with external applications.
    \item Establish a foundation for future multi-modal,
    multi-language and collaborative-user extensions.
\end{itemize}

\subsection{Measurable Performance Targets}
The system aims for the targets in Table~\ref{tab:targets}.

\begin{table}[!t]
\caption{Measurable performance targets.}
\label{tab:targets}
\centering
\renewcommand{\arraystretch}{1.2}
\begin{tabular}{@{}lll@{}}
\toprule
\textbf{Metric} & \textbf{Target} & \textbf{Method} \\
\midrule
Narrative coherence score   & $>$ 85\%  & Discourse analysis \\
User-preference alignment   & $>$ 80\%  & Likert vs. stated preference \\
User satisfaction rating    & $>$ 4.0/5 & Post-interaction survey \\
Context retention accuracy  & $>$ 90\%  & Cross-segment consistency \\
Response latency (cloud)    & $<$ 5\,s  & End-to-end timing \\
Session completion rate     & $>$ 75\%  & Sessions completed/started \\
Image generation success    & $>$ 95\%  & Successful first or retry \\
\bottomrule
\end{tabular}
\end{table}

\section{Scope of the Work}
To ensure focused and achievable outcomes, the following scope
boundaries are established:

\begin{itemize}
    \item \textbf{Text-first focus.} The system focuses on
    text-based narratives augmented with illustration, narration and
    ambient music; multi-modal generation of video is deferred to
    future work.
    \item \textbf{English language.} Initial development targets
    English-language generation; multi-language support is a
    secondary objective.
    \item \textbf{Fiction genres.} The system targets eight fictional
    narrative genres (fantasy, sci-fi, mystery, romance, thriller,
    adventure, horror, comedy).
    \item \textbf{Individual users.} The current scope addresses
    single-user interactions; multi-user collaborative storytelling
    is future work.
    \item \textbf{Short to medium length.} Target story lengths range
    from flash fiction (500 words) to short stories (5000 words),
    typically two to six chapters.
\end{itemize}

\section{Significance of the Study}
This work matters because it (i) introduces a deployable architectural
pattern --- the \texttt{custom\_elements} context channel --- that
generalises beyond storytelling to any LLM-driven application
requiring end-user-extensible personalisation; (ii) demonstrates that
free, hot-linked image and audio services can be combined into a
fluid multimodal experience when paced through a serial-throttling
queue; and (iii) provides a reproducible open-source codebase that
will serve as a baseline for future research on interactive narrative
generation.
